\uuid{IXqx}
\titre{Pgcd}
\auteur{bodin ; croizat ; sacre}
\organisation{exo7}
\licence{CC BY-NC-SA 4.0}
\datecreate{2022-03-08}
\niveau{L1}
\module{Algèbre}
\chapitre{Arithmétique dans Z}
\sousChapitre{Pgcd, ppcm, algorithme d'Euclide}
\difficulte{1}
\theme{pgcd}
\exo7qcmid{hilisit.arithmetique.003}

\contenu{
\begin{question}
\qclassification{103.03}
    Quelles sont les affirmations vraies ?
    \begin{answers}
        \good{$456$ est divisible par $3$.}
        \bad{$754$ est divisible par $4$.}
        \bad{$5552$ est divisible par $5$.}
        \bad{$987$ est divisible par $9$.}
    \end{answers}
    \begin{explanations}
     Critère de divisibilité par $3$ : la somme des chiffres est divisible par $3$. \\
     Critère de divisibilité par $2$ : le dernier chiffre est pair. \\
     Pour décider si un entier est divisible par $4$ : diviser l'entier par $2$ et appliquer le critère de divisibilité par $2$. \\
      Critère de divisibilité par $5$ : le dernier chiffre est $0$ ou $5$. \\
     Critère de divisibilité par $9$ : la somme des chiffres est divisible par $9$. \\
    \end{explanations}
\end{question}
}
